Encryption of clinical records at rest is usually specified as ``AES-256'' and judged by throughput, a specification that says nothing about whose record a ciphertext belongs to. An insider with write access to the database can move an encrypted field from one patient's row to another's, and the application, decrypting for the row it believes it is reading, presents the wrong patient's data with no error. We show this record-swap attack succeeds against shared-key AES-CBC with HMAC and against per-tenant AES-GCM without associated data, both common deployments. We present CPDPM, a construction that binds every ciphertext to its context of tenant, patient, resource type and field path: the context is encoded injectively, used to derive a per-context key from a versioned master key through HKDF, and bound again as associated data to a standard AEAD, with a key commitment against partitioning oracles, versioned rotation, and a single failure mode. For FHIR resources an encrypted manifest binds the set of encrypted fields, detecting deletion, reordering and per-field replay, and a streaming mode extends the guarantees to large payloads. Security reduces to standard AEAD notions. We also show by exhaustive check that the keyed byte-level pre-transform common in hybrid health-record schemes adds at most \KeyspaceLog{} bits and nothing when its key travels with the ciphertext. CPDPM rejects every manipulation the baselines accept, at a fixed cost of about \CpdpmFixedUs{} microseconds per operation over raw AES-GCM. Implementation, tests and experiments are released.
